\begin{center}
\textbf{Эффективные методы сканирования OCI-образов по метаданным}

\textit{Зиганшин Марк Радикович}
\end{center}

\vspace{0.7cm}

\begin{center}

В данной работе рассматривается задача повышения эффективности сканирования OCI-образов на наличие известных уязвимостей. Произведен обзор предметной области, изучено устройство OCI-образов, особенности их хранения в container registry, а также основные принципы, на которых базируется работа современных сканеров уязвимостей. Рассмотрены популярные сканеры Trivy, Clair и  Grype, для которых выявлен общий недостаток, связанный с избыточной загрузкой данных образа в процессе сканирования. В качестве решения предложено использовать подход, связанный с применением концепции ленивой загрузки образов. Разработан форк сканера Trivy, реализующий ленивую загрузку посредством формата eStargz. Разработан фреймворк для воспроизводимого бенчмарка, моделирующий типичный для container registry профиль нагрузки. По результатам экспериментальной оценки разработанный прототип демонстрирует кратное ускорение сканирования и существенное снижение объёма скачиваемых из registry данных по сравнению с оригинальным Trivy без потери точности

\end{center}

\newpage
